\paragraph{A retrieval-augmented baseline (answer to R5).}
To quantify how much of the observed citation hallucination is a grounding deficit rather than a generation behavior, we ran a minimal retrieval-augmented baseline for Claude Sonnet on the same 144 claims under the Baseline and Temporal conditions, injecting the top-5 Crossref records (by keyword, year-windowed under Temporal) into the prompt while holding the system prompt, output schema, temperature, and verification pipeline fixed.
Grounding raises Claude's Baseline existence rate from 0.381 to 0.699 ($\Delta = +0.318$, 95\% CI $[0.254,\,0.375]$).
Under the Temporal condition\,---\,the steepest closed-book drop (0.381$\rightarrow$0.119)\,---\,retrieval more than fully offsets the collapse: existence rises to 0.818 ($\Delta = +0.699$, 95\% CI $[0.649,\,0.746]$), exceeding even the unconstrained closed-book Baseline, while the temporal-violation rate stays at 0.035.
These two rows are added to Table~\ref{tab:delta}. Two caveats bound the interpretation: the injected candidates are by construction real Crossref records, so the gain partly reflects faithful reuse of supplied references rather than improved unaided recall, and 5 of the 288 RAG runs abstained from citing entirely (counted as zero existing). Even so, grounding eliminates most of the constraint-induced hallucination for this model, motivating the broader four-model, five-condition RAG study outlined in the now-shortened Future Work paragraph (Retrieval-augmented baselines).
